\documentclass[10pt]{article}

% Specify the margins and text size.
\setlength{\textwidth}{6.4in}
\setlength{\textheight}{9.5in}
\setlength{\oddsidemargin}{0pt}
\setlength{\evensidemargin}{0pt}
\setlength{\topmargin}{0pt}
\setlength{\hoffset}{.05in}
\setlength{\voffset}{-.8in}

\setlength{\parskip}{5pt}
\setlength{\parindent}{0pt}

% Load some fonts and symbol packages
\usepackage{latexsym}
\usepackage{pifont}       % contains 'star' symbol for counterinsurgency handbook title
\usepackage{yfonts} 
\usepackage{amsmath}
\usepackage{amsfonts}

\usepackage{graphicx}     % actually, this is loaded by pstricks
\usepackage[T1]{fontenc}
\usepackage{ifthen}
\usepackage{pstricks,pst-func,pst-text,pst-node,multido,pst-plot,calc,pst-3dplot}
%\usepackage[all]{xy}
%\usepackage{animate}

% The hyperref package inserts links into the pdf.
\definecolor{MyLinkColor}{rgb}{.1,.2,1}
\definecolor{MyCiteColor}{rgb}{.1,1,.2}
\definecolor{MyURLColor}{rgb}{.4,.4,.4}
\usepackage[backref=true,pagebackref=false,hyperindex,colorlinks=true,
  linkcolor=MyLinkColor,urlcolor=MyURLColor]{hyperref}


% The tweaklist package is something I found on the web.  It provides a simple interface
% for making changes to spacing used in the itemize and enumerate environments.  Comment
% this out if you don't care to use tweaklist.
\usepackage{tweaklist}
\renewcommand{\itemhook}{\setlength{\parskip}{2pt}\setlength{\parsep}%
{1pt}\setlength{\topsep}{0pt}\setlength{\itemsep}{0pt}}

\newcommand{\N}{\ensuremath\mathbb{N}}
\newcommand{\Q}{\ensuremath\mathbb{Q}}
\newcommand{\R}{\ensuremath\mathbb{R}}
\newcommand{\C}{\ensuremath\mathbb{C}}
\newcommand{\z}{\hphantom{0}}

\begin{document}
\pagestyle{empty}

\href{http://www.shepherd.edu}{\includegraphics[height=1.75cm]{logo-high-res.eps}}
\vspace{-1.69cm}

{\small
\begin{tabular}{cl}
\hspace{5in} & Math 314:  Statistics\\
                & Warmup Exercises\\
                & %6 February 2012
\end{tabular}
}
\medskip

{\setlength{\baselineskip}{1.05\baselineskip}

1. For the following data, (a)~write down the equation for the regression line, 
(c)~sketch the regression line and (c)~use the regression line to predict the $y$ value
if $x=130$.\vspace{-8pt}
\begin{center}
\begin{tabular}{llll}
x: & average $\approx 100$, & $\mbox{SD}\approx 15$\\
y: & average $\approx 110$, & $\mbox{SD}\approx 30$,  & $r\approx 0.8$\\
\end{tabular}
\end{center}
\vspace{1.25in}

2. For the following data, (a)~write down the equation for the regression line, 
(c)~sketch the regression line and (c)~use the regression line to predict the $y$ value
if $x=8$.\vspace{-8pt}
\begin{center}
\begin{tabular}{llll}
x: & average $\approx 10$, & $\mbox{SD}\approx 2$\\
y: & average $\approx 12$, & $\mbox{SD}\approx 6$,  & $r\approx -0.2$\\
\end{tabular}
\end{center}
%\vspace{1.0in}
\vfill\vfill

\href{http://www.shepherd.edu}{\includegraphics[height=1.75cm]{logo-high-res.eps}}
\vspace{-1.69cm}

{\small
\begin{tabular}{cl}
\hspace{5in} & Math 314:  Statistics\\
                & Warmup Exercises\\
                & %6 February 2012
\end{tabular}
}
\medskip

1. For the following data, (a)~write down the equation for the regression line, 
(c)~sketch the regression line and (c)~use the regression line to predict the $y$ value
if $x=130$.\vspace{-8pt}
\begin{center}
\begin{tabular}{llll}
x: & average $\approx 100$, & $\mbox{SD}\approx 15$\\
y: & average $\approx 110$, & $\mbox{SD}\approx 30$,  & $r\approx 0.8$\\
\end{tabular}
\end{center}
\vspace{1.25in}

2. For the following data, (a)~write down the equation for the regression line, 
(c)~sketch the regression line and (c)~use the regression line to predict the $y$ value
if $x=8$.\vspace{-8pt}
\begin{center}
\begin{tabular}{llll}
x: & average $\approx 10$, & $\mbox{SD}\approx 2$\\
y: & average $\approx 12$, & $\mbox{SD}\approx 6$,  & $r\approx -0.2$\\
\end{tabular}
\end{center}
\vfill
\eject
{\ }

3. For the small data set below left, $r\approx 0.76$.  Find $r$ for the other data sets.
\begin{center}
\begin{tabular}{ccccccc}
\begin{tabular}{cc}
$x$ & $y$\\[2pt]\hline
1 & 2\vphantom{\Large Y}\\
2 & 3\\
3 & 1\\
4 & 5\\
5 & 6
\end{tabular}
&\hspace{.5in} &
\begin{tabular}{cc}
$x$ & $y$\\[2pt]\hline
2 & 1\vphantom{\Large Y}\\
3 & 2 \\
1 & 3 \\
5 & 4 \\
6 & 5 
\end{tabular}
&\hspace{.5in} &
\begin{tabular}{cc}
$x$ & $y$\\[2pt]\hline
1 & 0\vphantom{\Large Y}\\
2 & 1\\
3 & -1\\
4 & 3\\
5 & 4
\end{tabular}
&\hspace{.5in} &
\begin{tabular}{cc}
$x$ & $y$\\[2pt]\hline
1 & 4\vphantom{\Large Y}\\
2 & 6\\
3 & 2\\
4 & 10\\
5 & 12
\end{tabular}
\end{tabular}
\end{center}
\vspace{.25in}

4. For the following data, calculate the correlation coefficient.
\begin{tabular}{|c|cccccc|}\hline
$x$ & 7 & 8 & 10 & 11 & 11 & 13\vphantom{\Large Y}\\
$y$ & 8 & 7 & 11 & 11 & 13 & 10\\\hline
\end{tabular}

\vfill


3. For the small data set below left, $r\approx 0.76$.  Find $r$ for the other data sets.
\begin{center}
\begin{tabular}{ccccccc}
\begin{tabular}{cc}
$x$ & $y$\\[2pt]\hline
1 & 2\vphantom{\Large Y}\\
2 & 3\\
3 & 1\\
4 & 5\\
5 & 6
\end{tabular}
&\hspace{.5in} &
\begin{tabular}{cc}
$x$ & $y$\\[2pt]\hline
2 & 1\vphantom{\Large Y}\\
3 & 2 \\
1 & 3 \\
5 & 4 \\
6 & 5 
\end{tabular}
&\hspace{.5in} &
\begin{tabular}{cc}
$x$ & $y$\\[2pt]\hline
1 & 0\vphantom{\Large Y}\\
2 & 1\\
3 & -1\\
4 & 3\\
5 & 4
\end{tabular}
&\hspace{.5in} &
\begin{tabular}{cc}
$x$ & $y$\\[2pt]\hline
1 & 4\vphantom{\Large Y}\\
2 & 6\\
3 & 2\\
4 & 10\\
5 & 12
\end{tabular}
\end{tabular}
\end{center}
\vspace{0.25in}

4. For the following data, calculate the correlation coefficient.
\begin{tabular}{|c|cccccc|}\hline
$x$ & 7 & 8 & 10 & 11 & 11 & 13\vphantom{\Large Y}\\
$y$ & 8 & 7 & 11 & 11 & 13 & 10\\\hline
\end{tabular}
\vspace{2in}

\eject
\end{document}
